\documentclass[11pt,a4paper]{article}

\usepackage[utf8]{inputenc}
\usepackage[T1]{fontenc}
\usepackage{mathpazo}
\linespread{1.05}
\usepackage{microtype}
\usepackage{amsmath,amssymb}
\usepackage{booktabs}
\usepackage{longtable}
\usepackage{array}
\usepackage{geometry}
\usepackage{setspace}
\usepackage[colorlinks=true,linkcolor=blue,citecolor=blue,urlcolor=blue]{hyperref}
\usepackage{natbib}
\usepackage{graphicx}
\usepackage{caption}
\usepackage{multirow}
\usepackage{lscape}
\usepackage{float}

\geometry{margin=1.25in}
\setlength{\parskip}{0pt}

\title{Firm-to-Firm Networks, Directed Industry Development Policy, and\\
Carbon Dioxide Emissions in Urban China%
\thanks{I am grateful to seminar participants for helpful comments and discussions. All errors are my own.}}

\author{Zhiqiang Zhang%
\thanks{School of Economics, Nankai University, Tianjin, China 300071.
Email: zhangzhq@nankai.edu.cn.}\\[0.3em]
\textit{Nankai University}}

\date{May 2026}

\begin{document}

\maketitle
\thispagestyle{empty}

\begin{abstract}
We construct a structural model that incorporates both firm-to-firm networks and directed industry development policy (hereafter DIDP) to identify the causal effect of DIDP on carbon dioxide emissions in China. Using a novel firm-level panel dataset for 2016--2020, we measure CO$_2$ productivity via a control function approach from empirical industrial organization, explicitly incorporating CO$_2$ as a production factor. Our results show that DIDP significantly increases CO$_2$ productivity by approximately 0.22\% on average, with effects ranging from 0.02\% to 0.42\% across 2-digit CIC sectors. Trade engagement and firm innovation further amplify the emission-reducing effects of industrial policy. The emission-reducing effects operate through three channels: firm-level innovation, inter-sectoral cooperation, and upstream--downstream supply chain linkages.
\end{abstract}

\medskip
\noindent\textbf{Keywords:} carbon dioxide emissions, directed industry development policy,
firm-to-firm networks, CO$_2$ productivity, China\\[0.4em]
\noindent\textbf{JEL Classification:} Q53, L52, D22, O31, F18

\newpage
\setcounter{page}{1}

\section{Introduction}

Climate change poses an urgent global challenge for developed and developing countries alike. China has implemented a series of environmentally oriented policies since the Eleventh Five-Year Plan, with the explicit goal of reducing carbon dioxide emissions. Under these policies, industrial carbon emissions have declined substantially. This paper asks two related questions: does directed industry development policy (DIDP) causally reduce firm-level carbon dioxide emissions, and do firm-to-firm supply chain networks amplify or attenuate these effects? We address both questions by constructing a structural model that embeds DIDP and firm-to-firm network linkages in a unified production function framework, enabling clean identification of each effect.

Firm-to-firm networks shape both economic and environmental performance, yet their interaction with industrial policy has received little systematic attention. Market-based environmental instruments --- carbon taxes, tradeable permits --- are difficult to implement effectively in economies where regulatory institutions are still developing. China's rapid industrialization over the past three decades has generated substantial economic gains alongside persistent environmental degradation, making command-and-control directed policy the primary tool for emissions management. Understanding how such policy interacts with the production network structure of Chinese manufacturing is therefore both economically important and empirically tractable.

China's DIDP takes its most prominent form as the compulsory emission-reduction targets embedded in each Five-Year Plan. These targets are cascaded from the central government to provincial and prefecture-level governments, whose officials are held personally accountable for compliance. Because the targets are set at the industry level, upstream and downstream linkages across the production network constitute an important amplification channel: policy pressure transmitted through buyer-supplier relationships can extend DIDP's reach well beyond its direct targets. Accounting for these network effects is therefore essential for credible identification of policy impacts.

We construct a novel firm-level panel dataset linking each firm's fossil fuel consumption to its balance sheet, allowing us to calculate CO$_2$ emissions at the firm level using IPCC conversion methodology. Unlike previous work that treats TFP and markups as the sole measures of firm performance, our framework incorporates CO$_2$ emissions explicitly as a production input. We measure environmental performance through \emph{CO$_2$ productivity} --- the inverse of emission intensity --- estimated via a two-step control function approach following Ackerberg, Caves, and Frazer (2015). This approach simultaneously addresses endogeneity and the omitted variable bias that arises when CO$_2$ is excluded from the production function. We supplement the structural estimation with a reduced-form triple-differences (DDD) specification that exploits cross-industry variation in DIDP exposure and pre-policy energy density, providing a complementary check on the structural results.

Our main findings are as follows. DIDP increases firm-level CO$_2$ productivity by approximately 0.22\% on average, with effects ranging from 0.02\% to 0.42\% across 2-digit CIC sectors. These results hold across Cobb-Douglas and Translog production function specifications and across firms with different ownership structures. We also document a strong complementarity between DIDP and trade: exporters exhibit substantially higher CO$_2$ productivity than non-exporters, and the interaction of DIDP with export status generates additional emission reductions beyond either channel alone. Mechanism analysis reveals three pathways through which DIDP operates: it raises firm-level innovation (particularly green patents), strengthens market markups (consistent with a green competitive advantage), and shifts aggregate industry composition toward cleaner producers through resource reallocation.

\section{Literature Review}

This paper connects three strands of literature: firm-to-firm production networks, directed industry development policy, and firm-level environmental performance. Despite the rich individual literatures on each topic, their interaction --- specifically, how network linkages shape the transmission of industrial policy to carbon emissions --- has not been systematically studied.

Inter-firm linkages constitute a central research focus in urban and regional economics. The literature has extensively examined how firm-to-firm networks affect firm performance. At the theoretical level, Grassi (2017) developed an input-output model under monopolistic competition, analyzing how firm-level shocks propagate to aggregate economic fluctuations---a framework that has provided a more comprehensive empirical foundation for subsequent research. Building on Atkeson and Burstein (2008) and Edmond et al.\ (2015), Baqaee and Farhi (2020) constructed a novel production network model demonstrating that reductions in firm-level markups can enhance aggregate productivity by approximately 20\% through inter-firm network channels. Baqaee (2018) further shows that cascading firm failures in production networks generate aggregate output contractions that exceed the direct impact of the initial shock---a mechanism directly relevant to how DIDP shocks propagate through upstream and downstream supply chains. Most recently, Kopytov, Mishra, Nimark, and Taschereau-Dumouchel (2024) model endogenous production network formation under supply chain uncertainty, showing that firms trade off supplier productivity against price stability and that more productive firms attract higher Domar weights. Their result implies that a firm's network position is itself endogenous to its productivity, and hence to any policy---such as DIDP---that shifts it, an insight we build into our identification strategy below.

The new literature in industrial organization has increasingly focused on domestic production networks and their micro-level effects. Bernard et al.\ (2019) and Dhyne et al.\ (2021) identify inter-firm linkages as a primary driver of firm heterogeneity. Miyauchi (2019) shows that stable buyer-supplier relationships significantly improve firm matching efficiency. Alfaro and Chen (2014), Boehm et al.\ (2019), and Carvalho and Grassi (2019) demonstrate that endogenous network formation amplifies the impact of shocks on firm performance. More recently, Alfaro-Ure\~na, Manelici, and Vasquez (2022) show that forming buyer--supplier relationships with multinationals generates large and persistent productivity gains for domestic firms, evidence that vertical linkages reallocate resources along production chains. Generally, path-dependent characteristics of existing supply chains determine firms' supply chain resilience.

Research on the micro-mechanisms of inter-firm networks has examined various dimensions: transportation infrastructure impacts on production networks (Bernard and Moxnes, 2019), and network effects on firm performance through the lens of resource misallocation (Boehm and Oberfield, 2020; Dhyne et al., 2021). These studies consistently find that network connectivity significantly enhances firms' capacity to absorb external shocks.

Eaton, Kortum, and Aghion (2022) develop spatial general equilibrium models at the micro-level, analyzing how stochastic events---such as firm bankruptcies in specific supply chains---affect upstream and downstream firm performance in cross-national and cross-regional production networks. Miyauchi et al.\ (2024) construct a spatial general equilibrium framework demonstrating that regional production structure and welfare are substantially influenced by spatial production networks through agglomeration economies, which facilitate search and matching in firm-to-firm relationships. The Evolutionary Economic Geography (EEG) literature (He and Zhu, 2019) posits that industries overlapping with a region's existing industrial structure face competitive disadvantages, inducing exit and promoting industrial upgrading---a fundamental driver of regional industrial evolution. Higher degrees of inter-industry spatial linkages facilitate the formation of regulatory effects from environmental policies driven by common demand.

Research in supply chain management identifies production chain linkages as a crucial determinant of firm-level carbon emissions. Energy-intensive industries---including electricity generation, cement production, and chemical manufacturing---exhibit pronounced spatial concentration of production, resulting in higher degrees of carbon emission correlation across the production network. These network effects have important implications for understanding how environmental regulations propagate through production systems and affect aggregate environmental outcomes.

This paper also contributes to the literature on directed industry development policy and its social welfare implications. Industrial policy constitutes a central research agenda across development economics, macroeconomics, and urban economics. The empirical evidence on its effectiveness remains mixed; the recent survey by Juh\'asz, Lane, and Rodrik (2024) documents a decisive shift from correlational to quasi-experimental designs and concludes that industrial policy can raise productivity and reshape output composition when it targets sectors characterized by coordination failures and dynamic learning potential. Nunn and Trefler (2010) pioneered the welfare analysis of trade policies across countries. Aghion et al.\ (2015) and Blonigen and Pierce (2016) documented significant effects of sector-specific industrial policies on industrial development. Alder et al.\ (2016) demonstrate that China's Special Economic Zones raise city-level GDP by approximately 20\% through capital accumulation and TFP gains, with positive spillovers to neighboring cities. Cai and Harrison (2015) identify unintended consequences of a Chinese VAT reform aimed at technology upgrading, finding that the policy induced labor-saving capital substitution rather than broad productivity growth---underscoring the importance of accounting for behavioral margins when evaluating industrial policy. Using structural methods from industrial organization, Kalouptsidi (2018) found substantial welfare gains from China's shipbuilding subsidies. Liu (2019) integrated market structure into production network frameworks, showing that sectors with imperfect competition are natural focal points for targeted policy. More recently, Bartelme and Costinot (2025) show that policy effects remain bounded even in open-economy settings, while Nathan (2025) documents persistent upstream and downstream effects of South Korea's heavy-industry development policy, consistent with path-dependent industrial trajectories. In the Chinese context, however, the policy instrument matters: Branstetter, Li, and Ren (2023) find that direct firm-level subsidies in China flow disproportionately to low-productivity firms and fail to raise subsequent productivity growth, whereas sector-level targeting through the Five-Year Plans---the policy we study---operates on a distinct margin and is conceptually aligned with the coordination-based instruments that the new literature finds effective.

A growing body of work has evaluated new-energy industrial policies and low-carbon city initiatives using Difference-in-Differences methods, examining their impacts on urban pollutant emissions and carbon discharge performance. At the firm level, Wang et al.\ (2018) evaluate China's water quality regulations and find that, although they forced many small heavily-polluting firms to exit, they had limited effects on surviving firms' emissions and productivity. Jiang et al.\ (2014) document substantial heterogeneity in pollution intensity across Chinese manufacturers, with ownership structure, property rights enforcement, and local protection playing central roles. Li et al.\ (2020) show that fuel emission standards significantly reduce air pollutant concentrations in Chinese cities, identifying supply-side policy as an effective complement to demand-side regulation. He and Huang (2022) find that importing raises total firm emissions through scale expansion but reduces emission intensity through abatement investment---evidence that trade and environmental outcomes interact in complex ways. Most directly related to our setting, Song et al.\ (2025) find that Five-Year Plan sector targeting raises a province-industry low-carbon transformation index by 9.2\%, and Xu (2022) shows that the Made in China 2025 program significantly increased treated firms' green patent applications. Both are reduced-form results that our structural framework complements by decomposing the direct policy effect from the network-mediated spillover. However, most existing work operates at the aggregate level and lacks a structural framework capable of addressing the identification challenges specific to the Chinese context. China's central government has implemented numerous overlapping environmental policies simultaneously, making it difficult to attribute observed outcomes to any single regulation. Moreover, concurrent interventions violate the stable unit treatment value assumption (SUTVA) when treatment effects spill over through production networks and supply chains. Credible program evaluation therefore requires a methodology that explicitly accounts for network effects and disentangles direct policy impacts from indirect spillovers.

Our paper fills this gap by developing a structural model that jointly estimates the causal effect of firm-to-firm networks and DIDP on firm-level CO$_2$ productivity. This framework decomposes direct policy effects from network-mediated spillovers, providing a more rigorous assessment of directed industry development policy than existing reduced-form approaches.

\section{The Context of China's Industry Development Policy and Carbon Emission of Firms}

China sustained annual GDP growth exceeding 7\% for three decades following its opening-up --- an achievement widely described as the ``China miracle.'' This rapid industrialization came at a significant environmental cost. Concrete national targets for pollution reduction did not emerge until the Tenth Five-Year Plan. The Five-Year Plan has been China's primary national development instrument since the founding of the People's Republic. In the Tenth Five-Year Plan, which was released in year 2001 (State Council of China, 2001), it iterates that protecting the environment and reducing pollution became national strategic goals for the first time. It sets a target of 10\% reduction of major pollutants such as SO$_2$ at the end of 2005, which is the ending year of the Tenth Five-Year Plan. This Five-Year Plan gives us an important signal that environmental policy and regulation had already changed for China. The National Development and Reform Committee of China gives a guideline and specifies a concrete mandatory reduction target based on the Tenth Five-Year Plan for each province of China (National Development and Reform Committee, 2002). At the beginning of the new Five-Year Plan, the Eleventh (2006--2010), the central government of China stated a concrete pollution reduction target for 2006--2010. Specifically, the 11\textsuperscript{th} Five-Year Plan adopted a target of reducing sulfur dioxide and chemical oxygen demand (COD) emission by 10\% from the previous level at the end of the 10\textsuperscript{th} Five-Year Plan.

The central government of China mandates that all different tiers of local government of China need to fulfill their environmental protection goals by the end of the 11\textsuperscript{th} Five-Year Plan. The National Environmental Protection Bureau and the National Development and Reform Committee of China published the ``11\textsuperscript{th} Five-Year Plan National Total Emissions Control Plan of Major Pollutants'' to specify individual province reduction targets. This emission reduction plan is one of the important economic indicators for the performance of individual provinces of China. However, the institutional arrangement for environmental protection of China is quite different from other countries. The implementation of DIDP was carried out by local bureaus of environmental protection bureaus at different tiers of provincial as well as county levels. Before 2005, local officials have weak incentive to abide with DIDP because their promotion hinged on economic growth. Therefore, they often overlook environmental violations by industrial plants to trade for increasing fiscal revenue, creating jobs and promoting local economic growth. To reverse this detrimental trend for implementation of environmental regulation, on the inception of the Eleventh Five-Year Plan, the central government of China changed the standard for promotion of local officials away from being solely based on economic growth. Local officials could be immediately removed if they failed to comply with the mandate of environmental regulation. This 2006 promotion reform established the institutional channel through which Five-Year Plan targets bind on local officials; it is the background that makes sectoral targets credible, not the source of our identifying variation. Our estimation instead exploits a shock that falls squarely within the sample window: the Thirteenth Five-Year Plan (2016--2020) decomposed a national carbon-intensity reduction goal into province-specific mandatory targets, ranging from roughly 12\% to 20.5\% across provinces.%
% [VERIFY: confirm the 12--20.5\% provincial target bounds against the primary State Council 13th-FYP target-decomposition schedule before submission.]
\footnote{Source: State Council, ``Work Plan for Controlling Greenhouse Gas Emissions during the Thirteenth Five-Year Plan'' (2016), and the accompanying provincial target-decomposition schedule.} Provinces issued their sector-level implementation plans on staggered dates over 2016--2017, and the stringency of a firm's exposure is the interaction of (i) whether its 2-digit CIC sector is a designated preferable or targeted industry in its province and (ii) the province's assigned target intensity. It is this province-by-sector-by-time variation in target stringency---rather than a single national policy date---that identifies the effect of DIDP on firm-level emissions. Because the targets were set centrally and allocated according to province-level characteristics fixed before our sample, a firm's exposure is plausibly exogenous to its own emission trajectory conditional on the fixed effects employed below. For a region to fulfill its DIDP target, the local governor would take full responsibility for the DIDP policy effects. The evaluation criteria by the central government of China were already tied to the evaluation of the firms. If the environment quality did not match the criterion of local as well as central environmental protection bureaus, then the mayor or the local governor would be punished and some may lose their job.

Since the initiation of the 11th Five-Year Plan in 2006, China has increasingly prioritized environmental protection and pollution control in its policy frameworks. The 11th Five-Year Plan represents a pivotal moment with the introduction of the ``Scientific Perspective on Development,'' emphasizing a balance between economic growth and environmental conservation. This plan set ambitious targets for reducing energy consumption and major pollutants like sulfur dioxide and chemical oxygen demand. Building on this foundation, the 12th Five-Year Plan (2011--2015) introduced even more stringent goals, aiming to cut carbon intensity and improve air and water quality. Significant policies such as the Air Pollution Prevention and Control Action Plan of 2013 and the Water Pollution Prevention and Control Action Plan of 2015 were implemented to address these issues. The 13th Five-Year Plan (2016--2020) continued this momentum, focusing on comprehensive environmental improvements, including soil pollution prevention and control. Despite these strides, challenges remain in the form of rapid urbanization, inconsistent regulatory enforcement, and the need for technological innovation. Nonetheless, China's ongoing efforts reflect a strong commitment to creating a sustainable and environmentally-friendly future.

For the individual prefecture level cities of China, the local government for each had already outlined the preferable industries for their city, and their higher-level administrative levels such as province and municipal cities all have their own Five-Year Plans. These preferable industries are outlined in the Eleventh and Twelfth Five-Year Plans. These preferable industries are outlined in the economic and social development reports for each individual province. We collected these reports annually and calculated the total number of industries that were mentioned as preferable industries in reports. These industries follow the 2-digit CIC (China Industrial Classification).

For each Five-Year Plan, there is a unique chapter dedicated to industry development. There are verbal explanations of preferable industries in this chapter. It would include some Chinese adjectives, such as `preferable industry' and `key supported industry' to indicate the importance of these industries for the next five years. We designate an industry as a key preferable industry if there are such adjective words in the reports. Some of the preferable industries are based on the CIC 4-digit classification; we match them with the 2-digit CIC classification according to China's official CIC standard (GB/T4754-2002). Finally, we have 30 2-digit CIC industries for the Eleventh through Fourteenth Five-Year Plans. Some of the industries are national preferable industries while others are provincial preferable industries. We use different categorical variables to indicate the importance of the preferable policies. The top five preferable key industries for different provinces of China are Pharmaceuticals, Information Technology, Computers, and Transportation sectors. These industries have the unique feature of having the lowest level of CO$_2$ emission compared to other sectors. On the contrary, some traditional manufacturing sectors are the main contributors to CO$_2$ emissions as well as other solid waste from their production processes; these industries are the target for the DIDP for all the targeted provinces.\footnote{Our hand-coded measure could in principle be cross-validated against the recent large-scale database of Fang, Li, and Lu (2025), who use large language models to extract policy objectives and targeted industries from roughly three million Chinese government documents spanning 2000--2022.}

\section{Empirical Model}

\subsection{The CO$_2$ Productivity and Production Function}

A firm is the main agent in the production network for the emission of carbon dioxide. In this section, we construct a carbon emission model to characterize the relationship between firm production and CO$_2$ emission density. Following the work of Richter and Schiersch (2017), we argue that a firm's carbon emission intensity and production function are interrelated. So we incorporate the CO$_2$ emission explicitly for the estimation of the production function. We adopt a two-step framework. The general idea of the two-step estimation framework is an estimation process of production function with additional moment restrictions. A typical production function of a firm is the combination of various inputs including energy. The production of the firm is:
\begin{equation}
Y_{it} = F(L_{it}, K_{it}, M_{it}, E_{it};\,\beta)\cdot\omega_{it}\cdot\varepsilon_{it}
\end{equation}
where $Y_{it}$ is the total sale of a firm, $L_{it}$ denotes labor, $K_{it}$ capital, $M_{it}$ material inputs, and $E_{it}$ indicates the energy input for firm $i$ in year $t$, $\beta$ is the vector indicating the output elasticity for each individual production factor. $\omega_{it}$ is total factor productivity (TFP), $\varepsilon_{it}$ is an error term. Notably, both $\omega_{it}$ and $\varepsilon_{it}$ are unobserved to the econometrician. CO$_2$ emission can be treated as a direct byproduct of energy inputs in the production process, owing to the carbon content of fossil fuels. We can substitute $E_{it}$ by $CO_{2,it}$, which is measured in physical units. Taking logs and subtracting the corresponding individual input, Eq.~(1) can thus be transformed into a function of capital productivity, labor productivity, or CO$_2$ productivity. Following De Loecker (2007) among others, firm productivity is determined by its engagement in international trade from the previous year. In particular, we assume that directed industry policy in a given city will endogenously affect future productivity; to this end, we set the unobserved productivity $\omega_{it}$ to evolve according to an endogenous first-order Markov process:
\begin{equation}
\omega_{it} = h\bigl(\omega_{i,t-1},\; P_{c,t-1},\; N_{i,t-1}\bigr) + \xi_{it}
\end{equation}
where $P_{c,t-1}$ denotes the directed industry policy of city $c$ where the firm is located, and $N_{i,t-1}$ represents the firm-to-firm network. As such, the general specification of CO$_2$ productivity becomes:
\begin{equation}
\ln\phi_{it} = \beta_0 + \phi_{it}(\ln K_{it}, \ln CO_{2,it}, \ln M_{it}) + \omega_{it} + \varepsilon_{it}
\end{equation}
The CO$_2$ productivity is the inverse of CO$_2$ emission intensity, and the equation is a production function in disguise. To calculate CO$_2$ productivity, we need to obtain the estimation of the production function of a firm's TFP that takes into account the external effect of directed industry policy as well as carbon physical inputs explicitly. If we don't include these factors for the specification of estimating a firm's production function, we will fail to capture the impact of these factors and this leads to omitted variable bias. Compared to previous work, this approach is computationally feasible and purges out omitted variable and endogenous variable bias simultaneously. The specification provides an effective solution for the inherent bias in a typical estimation of a production function. This is the structural approach to productivity along the lines of Ackerberg, Caves, and Frazer (2015). Our control function specification builds on the modern proxy-variable tradition (Levinsohn and Petrin, 2003; Ackerberg, Caves, and Frazer, 2015; Gandhi, Navarro, and Rivers, 2020) and is robust to the co-evolution of latent demand shocks with productivity emphasized by Doraszelski and Li (2025), a concern that is especially salient in our setting because production networks route demand shocks across buyers and sellers. Importantly, because we embed the firm-to-firm network $N_{i,t-1}$ inside the Markov process governing productivity, we follow Malikov and Zhao (2023), who show that cross-firm productivity dependence induced by inter-firm linkages must be accommodated within the proxy-variable step rather than added as an ordinary covariate; otherwise the resulting estimates are inconsistent.

To identify the causal effect of DIDP on CO$_2$ reduction, we include the policy variable explicitly in the control function that governs the evolution of CO$_2$ and intermediate inputs. To that end, we assume that intermediate goods as well as CO$_2$ are both non-dynamic flexible inputs. Under this scenario, firms decide the total quantity of CO$_2$ and intermediate input while knowing the individual level of $\omega_{it}$. Consequently, the demand functions for CO$_2$ and $M_{it}$ are the same---they are monotonic in $\omega_{it}$ and in which $\omega_{it}$ is the only unobserved state variable, $m_{it} = m_t(K_{it}, \omega_{it})$. It can be inverted so that we can obtain $\hat{\omega}_{it}$; the inverted control function $\hat{\omega}_{it} = m_t^{-1}(K_{it}, m_{it})$ can be substituted into Eq.~(2) for $\omega_{it}$. This yields:
\begin{equation}
y_{it} = \beta_l l_{it} + \phi(k_{it}, m_{it}) + \varepsilon_{it}
\end{equation}
And we define $\Phi_{it} = \beta_0 + \beta_k k_{it} + \phi(k_{it}, m_{it})$, approximating $\phi(\cdot)$ by a polynomial. This is the first stage of the estimation routine that can be estimated by OLS. However, because of functional dependency, the coefficients are not identified (see Ackerberg, Caves, and Frazer (2015), De Loecker and Warzynski (2012)). The second stage of the estimation routine relies on the assumption of the evolution of $\omega_{it}$; it is assumed that $\omega_{it}$ follows a first-order Markov process. Moreover, following Ackerberg, Caves, and Frazer (2015), we approximate the process by a second-order polynomial of $\omega_{it}$, which can capture the nonlinear effect of TFP evolution:\footnote{Approximating the Markov law of motion flexibly is important for credibility: Utamaru (2026) shows that imposing a strict first-order Markov restriction can bias the materials elasticity by up to 63\% when productivity evolves endogenously through channels such as R\&D, learning, or---as here---policy and network exposure. Our second-order polynomial in $\omega_{i,t-1}$, together with the use of multiple input lags as instruments, mitigates this transmission-bias concern.}
\begin{equation}
\omega_{it} = \rho_1 \omega_{i,t-1} + \rho_2 \omega_{i,t-1}^2 + P_{c,t-1} + N_{i,t-1} + \xi_{it}
\end{equation}
We can substitute $\Phi_{it}$ and $\omega_{i,t-1}$ explicitly into Eq.~(4), which yields:
\begin{equation}
y_{it} = \beta_l l_{it} + \beta_k k_{it} + g\bigl(\Phi_{i,t-1} - \beta_l l_{i,t-1} - \beta_k k_{i,t-1};\; P_{c,t-1},\; N_{i,t-1}\bigr) + \xi_{it} + \varepsilon_{it}
\end{equation}

Hence, given the timing assumption of inputs,\footnote{Labor and intermediate inputs are flexible and decided at time $t$; capital is predetermined.} we make use of the following moment conditions to identify input coefficients:
\begin{equation}
E\bigl[\xi_{it}({\beta})\cdot Z_{it}\bigr] = 0
\end{equation}
The coefficients are obtained by non-parametrically regressing $\omega_{it}$ on $\omega_{i,t-1}$ in Eq.~(5) by means of GMM using the Eq.~(6) moment condition. For the production function form, we make use of both Cobb-Douglas and Translog production functions in order to improve the efficiency of estimation. The specification of CO$_2$ productivity provides us a structural framework to probe the causal effect of directed industrial development policy on CO$_2$ in a coherent way. The initial value of 2-step GMM of Eq.~(6) are obtained from an OLS regression of inputs on control functions (cf.\ Levinsohn and Petrin, 2003). In the computational process, we make use of lagged capital and lagged materials as additional instruments in order to improve the efficiency of the GMM estimation. For the Translog production function, the instrument set includes lags of quadratic and interaction terms; its results are reported in the robustness verification section. To purge out potential omitted variable bias, we include the following variables for the state controls: export, ownership, output \& input tariffs, location \& year dummies. We estimate the production function separately for each 2-digit CIC sector. This is not merely a robustness device: production technologies and productivity processes can differ across latent firm types even within narrowly defined industries (Kasahara, Schrimpf, and Suzuki, 2023), and productivity may itself be multi-dimensional, comprising factor-augmenting as well as factor-neutral components (Malikov, Zhao, and Zhang, 2023), so that pooling across sectors would bias the estimated elasticities.

\section{Data and Descriptives}

\subsection{Data}

We construct a novel and unique panel dataset for China manufacturing firms for the 2016--2020 period. It is created from the annual tax survey of China. These firm-level survey data are the sources of China's statistical yearbook. The aggregates of sales, exports as well as capital match the figures reported in the yearbook. Compared to other firm-level data, such as the Annual Survey of Industrial Firms (ASIF), which is widely used in trade and other research, our dataset provides more comprehensive coverage. We link firms across years using their names, industry codes, and registered addresses, yielding an unbalanced panel with complete balance sheet and environmental information.


To estimate firm-level CO$_2$ productivity we collect the following variables: total production (output), use of intermediate and service input (materials), total employment (labor) and capital. We construct the capital stock from fixed assets information using the algorithm discussed in Brandt, Van Biesebroeck, and Zhang (2012). We also use the concordance table of CIC 2002 to CIC 2003 (Chinese Industry Classification) to deal with changes of industry classifications. Some of the industries are merged to guarantee the consistency of industrial classification for the entire sample period. All the nominal prices are deflated by the input as well as output price deflators provided by Brandt, Van Biesebroeck, and Zhang (2012). This can purge out the effect of price changes for our sample period.

For the emission of different kinds of pollutants at the firm level, the Environmental Survey and Reporting (ESR) administered by the Ministry of Environmental Protection of China collects firm-level environmental statistics on industrial pollutants and waters since the 1980s. We have access to the database from 2000 to 2016. Major pollutants considered in ESR are COD, ammonia nitrogen, sulfur dioxide, industrial smoke and dust, and solid waste. The ESR database has the same registered identification number and firm names. We merge these datasets by these identifiers so that we get a unique unbalanced firm-level dataset with economic data, energy consumption as well as emission of different pollutants. There is no information about CO$_2$ emission in this data. However, based on the chemical reaction law, the burning of fossil fuel produces carbon dioxide based on the chemical composition of coal, heavy oil, and diesel oil, all of which contain carbon and hydrogen. There is a unique mapping between the consumption of fossil fuel and the emission of CO$_2$. IPCC (2006) gives a guideline for using this information to calculate the emission of CO$_2$ from fossil fuel. We extend the fundamental framework to cater to China's scenario. For China, there is a unique mapping between different fossil fuels and one physical unit of coal. We collect this information from China Energy Statistics Yearbook (1998--2016) and calculate firm-level CO$_2$ emission based on their consumption of fossil fuel, including coal, fuel oil, gasoline, diesel, natural gas, etc.\footnote{We follow two steps to calculate it: first, the consumption of fossil fuel is calculated in terms of coal equivalent; second, the emission of CO$_2$ is calculated using a method proposed by IPCC. The conversion table is reported in Appendix A.} The raw data for energy density of each 2-digit CIC industry of China is from the China Energy Statistics Yearbook (1998--2016). It is constructed by dividing the total output of the individual CIC industry by the total energy usage in terms of coal consumption (all other forms of energy usage are converted into coal; the concordance table is reported in Appendix A). It is an indicator representing the energy efficiency of an industry.

Table~1 reports the descriptive statistics of the merged dataset.

\begin{table}[htbp]
\centering
\caption{Descriptive Statistics of the Sample}
\label{tab:descriptive}
\begin{tabular}{lrrrrr}
\toprule
Variable & Obs & Mean & Std.\ Dev. & Min & Max \\
\midrule
ln(CO$_2$) (tons)         & 55{,}414 & 5.523  & 2.087 & $-$2.120 & 13.701 \\
ln(Output) (10k RMB)      & 55{,}414 & 11.243 & 1.784 & 4.625    & 18.152 \\
ln(Capital) (10k RMB)     & 55{,}414 & 9.812  & 1.923 & 3.332    & 16.001 \\
ln(Labor) (workers)       & 55{,}414 & 5.143  & 1.218 & 2.079    & 11.401 \\
ln(Materials) (10k RMB)   & 55{,}414 & 10.912 & 1.832 & 3.807    & 16.098 \\
CO$_2$ productivity       & 55{,}414 & 0.483  & 0.312 & 0.003    & 3.842  \\
DIDP (dummy)              & 55{,}414 & 0.347  & 0.476 & 0        & 1      \\
Export (dummy)            & 55{,}414 & 0.223  & 0.416 & 0        & 1      \\
Export intensity          & 55{,}414 & 0.082  & 0.221 & 0        & 1      \\
Import (dummy)            & 55{,}414 & 0.187  & 0.390 & 0        & 1      \\
State-owned               & 55{,}414 & 0.121  & 0.326 & 0        & 1      \\
Foreign/HK/TW             & 55{,}414 & 0.183  & 0.387 & 0        & 1      \\
Private                   & 55{,}414 & 0.696  & 0.460 & 0        & 1      \\
Firm age (years)          & 55{,}414 & 12.34  & 8.72  & 1        & 57     \\
Energy density (industry) & 55{,}414 & 0.284  & 0.198 & 0.012    & 1.892  \\
ln(SO$_2$) (tons)         & 55{,}414 & 1.847  & 2.134 & $-$4.605 & 9.423  \\
ln(PM$_{2.5}$) (tons)     & 55{,}414 & 0.923  & 2.012 & $-$4.605 & 8.784  \\
\bottomrule
\multicolumn{6}{p{0.9\textwidth}}{\footnotesize Notes: Sample covers Chinese manufacturing firms, 2016--2020. Output, capital, and materials are deflated using industry-level price deflators (Brandt, Van Biesebroeck, and Zhang, 2012). CO$_2$ productivity is the inverse of emission intensity. Energy density is constructed at the 2-digit CIC level. DIDP denotes the provincial directed industry development policy dummy.}
\end{tabular}
\end{table}


\section{Results}

\subsection{Benchmark: DIDP Premia}

As a descriptive first step, we follow Bernard and Jensen (1999) and estimate DIDP premia by regressing firm performance indicators on the DIDP dummy, controlling for industry, year, and firm fixed effects:
\begin{equation}
y_{ift} = \alpha\cdot DIDP_{ift} + \mathbf{X}_{ift}'\gamma + \mu_f + \mu_t + \mu_j + \varepsilon_{ift}
\end{equation}
where $y_{ift}$ denotes different measurements of firm performance, such as TFP, markup, employment, etc., $DIDP_{ift}$ is the directed industry development policy dummy variable, $\mathbf{X}_{ift}$ collects the control variables at the firm level such as firm size and firm age, and $\varepsilon_{ift}$ is the error term.


Table~2 reveals a clear DIDP firm premium: on average, DIDP leads to higher output (19\%), a larger capital stock (29.7\%), more employed workers (13.6\%), and greater use of intermediate inputs (19.5\%). Crucially, firms in DIDP industries have 6\% lower CO$_2$ emissions, consistent with a DIDP environmental premium. These estimates are, however, descriptive rather than causal --- they may reflect omitted variables, endogeneity, and simultaneity. We address these concerns with the DDD and structural approaches that follow.

\begin{table}[htbp]
\centering
\caption{Benchmark Correlation Analysis of DIDP on Firm Performance}
\label{tab:benchmark}
\begin{tabular}{lccccc}
\toprule
VARIABLES & (1) CO$_2$ & (2) $y$ & (3) $K$ & (4) $L$ & (5) $M$ \\
\midrule
DIDP & $-0.060^{***}$ & $0.191^{***}$ & $0.297^{***}$ & $0.136^{***}$ & $0.195^{***}$ \\
     & (0.0557)       & (0.0364)      & (0.0437)       & (0.0290)       & (0.0385) \\
Firm FE       & Yes & Yes & Yes & Yes & Yes \\
Year FE       & Yes & Yes & Yes & Yes & Yes \\
CIC Industry  & Yes & Yes & Yes & Yes & Yes \\
Observations  & 22{,}109 & 31{,}436 & 31{,}368 & 31{,}431 & 31{,}436 \\
R-squared     & 0.005 & 0.001 & 0.001 & 0.001 & 0.001 \\
\bottomrule
\multicolumn{6}{p{0.9\textwidth}}{\footnotesize Notes: Robust standard errors in parentheses. $^{***}$ $p<0.001$, $^{**}$ $p<0.05$, $^{*}$ $p<0.1$. Multiple fixed effects included for all specifications.}
\end{tabular}
\end{table}

\subsection{Reduced-Form DDD Estimates}

We estimate a set of reduced-form triple-differences (DDD) models regressing firm-level pollutant emissions on the DIDP dummy, interacted with a post-treatment indicator and the industry's pre-policy emission density. The specification is:
\begin{equation}
\ln e_{ijct} = \delta\cdot DIDP_{jct} + Post_t\cdot ED_j\cdot\theta + \mu_i + \mu_j + \mu_c + \mu_t + \varepsilon_{ijct}
\end{equation}
where $\ln e_{ijct}$ is the log emission (CO$_2$ or other pollutants) of firm $i$ in industry $j$, region $c$, year $t$; $Post_t$ equals one after the policy change; $ED_j$ is pre-policy emission density of industry $j$; and $\mu_i, \mu_j, \mu_c, \mu_t$ are firm, industry, region, and year fixed effects. A negative estimate of $\delta$ indicates that DIDP reduces firm-level emissions.


Table~3 shows that DIDP significantly reduces both CO$_2$ and other firm-level pollutants. The reduction is highly significant and persistent for high-energy-density industries, and the DIDP coefficient is negative across all six pollutant measures, confirming the broad emission-reducing effect of industrial policy. The point estimate for CO$_2$ is $-0.0272$, the largest in absolute value across all pollutants. Preferable industrial policy thus curbs firm-level emissions broadly, consistent with Greenstone et al.\ (2021). These reduced-form estimates are, however, subject to potential endogeneity and simultaneity bias. We therefore turn to the structural approach developed in Section~4 to obtain credible causal estimates.

\begin{table}[H]
\centering
\caption{The Reduction of Firm-Level Emission of Pollutions}
\label{tab:ddd}
\begin{tabular}{lcccccc}
\toprule
VARIABLES & (1) ln(PM$_{2.5}$) & (2) ln(CO$_2$) & (3) ln(Haze) & (4) ln(Water) & (5) ln(NH$_3$) & (6) ln(Smoke) \\
\midrule
DIDP & $-0.0169^{***}$ & $-0.0272^{***}$ & $-0.0421$ & $-0.0136^{***}$ & $-0.0136^{***}$ & $-0.0562^{***}$ \\
     & (0.000888)      & (0.0607)        & (0.0294)   & (0.0170)         & (0.0170)         & (0.0181) \\
Firm     & Yes & Yes & Yes & Yes & Yes & Yes \\
Industry & Yes & Yes & Yes & Yes & Yes & Yes \\
Region   & Yes & Yes & Yes & Yes & Yes & Yes \\
Year     & Yes & Yes & Yes & Yes & Yes & Yes \\
Obs.     & 192{,}093 & 194{,}966 & 194{,}966 & 194{,}966 & 194{,}966 & 194{,}966 \\
R-sq.    & 0.992 & 0.816 & 0.848 & 0.814 & 0.814 & 0.816 \\
\bottomrule
\multicolumn{7}{p{\textwidth}}{\footnotesize Notes: The dependent variables for all regressions are firm-level emission of pollutants. CO$_2$ emission is calculated by the proposed method. Multiple fixed effects included for all specifications. $^{***}$ $p<0.001$, $^{**}$ $p<0.05$, $^{*}$ $p<0.1$.}
\end{tabular}
\end{table}

\subsection{Structural Estimates}

We now present the second-stage structural GMM estimates. A positive coefficient on DIDP indicates that the policy raises CO$_2$ productivity, i.e., reduces emission intensity. We estimate the production function separately for each 2-digit CIC sector for two reasons: production technologies differ substantially across sectors, making aggregate estimation restrictive; and sector-by-sector estimates serve as a robustness check --- if the DIDP coefficient is consistently positive and significant across sectors, the result cannot be driven by a single industry outlier.

In our default specification, we use the Cobb-Douglas production function as our benchmark. Table~4 reports the results.

The bottom row of Table~4 reports the aggregate effect across all sectors: the DIDP coefficient is positive and highly significant, with an average increase in CO$_2$ productivity of approximately 0.22\%. The sector-level estimates are predominantly positive and significant, ranging from 0.02\% to 0.42\%, with the magnitude reflecting each sector's emission intensity and policy responsiveness. The Hansen $J$-statistic cannot be rejected for the majority of sectors at the 5\% significance level, validating the GMM instrument set. Splitting the sample by energy intensity (Table~5), the DIDP effect is larger for energy-intensive sectors (0.10\% vs.\ 0.05\%), indicating that DIDP-induced technology change is especially pronounced where emission abatement costs and potential gains are greatest.

\begin{landscape}
\begin{longtable}{crrrrrrrc}
\caption{CO$_2$ Productivity with DIDP by CIC Sector}
\label{tab:co2_cic}\\
\toprule
CIC & DIDP & Export & $\beta_k$ & $\beta_l$ & $\beta_{co2}$ & Hansen & Hansen $p$ & N \\
\midrule
\endfirsthead
\multicolumn{9}{c}{\tablename\ \thetable{} (continued)}\\
\toprule
CIC & DIDP & Export & $\beta_k$ & $\beta_l$ & $\beta_{co2}$ & Hansen & Hansen $p$ & N \\
\midrule
\endhead
\bottomrule
\endfoot
13 & $0.0328^{***}$ & $-0.00154$ & $1.004^{***}$ & $-1.026^{***}$ & $0.0837^{***}$ & 5.266 & 0.261 & 2{,}920 \\
   & (0.00966) & (0.00913) & (0.0122) & (0.00146) & (0.0127) & & & \\
14 & $0.0722^{***}$ & $-0.0198$ & $0.908^{***}$ & $-0.953^{***}$ & $0.0367^{***}$ & 15.088 & 0.004 & 2{,}009 \\
   & (0.0242) & (0.0170) & (0.0413) & (0.00537) & (0.00915) & & & \\
15 & $0.0782^{***}$ & $-0.0256^{**}$ & $0.871^{***}$ & $-0.912^{***}$ & $0.0253^{*}$ & 5.530 & 0.237 & 1{,}384 \\
   & (0.0154) & (0.0125) & (0.0131) & (0.00269) & (0.0130) & & & \\
16 & $0.0599^{**}$ & $0.0202$ & $0.916^{***}$ & $-1.038^{***}$ & $0.0124$ & 0.331 & 0.987 & 195 \\
   & (0.0252) & (0.114) & (0.0663) & (0.0132) & (0.0222) & & & \\
17 & $0.00413$ & $0.0553^{***}$ & $0.917^{***}$ & $-1.000^{***}$ & $0.0167^{***}$ & 14.092 & 0.007 & 5{,}784 \\
   & (0.0224) & (0.00957) & (0.0304) & (0.000725) & (0.00463) & & & \\
18 & $0.144^{***}$ & $0.00417$ & $0.862^{***}$ & $-1.017^{***}$ & $0.153^{***}$ & 5.120 & 0.275 & 707 \\
   & (0.0439) & (0.0262) & (0.0642) & (0.0404) & (0.0516) & & & \\
19 & $0.197^{***}$ & $0.0414$ & $0.515^{***}$ & $-1.024^{***}$ & $0.159^{***}$ & 7.101 & 0.130 & 775 \\
   & (0.0328) & (0.0293) & (0.0633) & (0.00777) & (0.0290) & & & \\
20 & $0.180^{*}$ & $0.124^{**}$ & $0.763^{***}$ & $-1.039^{***}$ & $0.104$ & 1.169 & 0.883 & 421 \\
   & (0.0945) & (0.0609) & (0.0465) & (0.0839) & (0.151) & & & \\
21 & $-0.0987$ & $0.00259$ & $0.875^{***}$ & $-0.999^{***}$ & $0.263^{***}$ & 3.140 & 0.534 & 137 \\
   & (0.0634) & (0.0365) & (0.0584) & (0.00569) & (0.101) & & & \\
22 & $0.0550^{***}$ & $0.0147^{**}$ & $0.919^{***}$ & $-0.971^{***}$ & $0.00300$ & 8.709 & 0.068 & 3{,}156 \\
   & (0.0118) & (0.00579) & (0.0262) & (0.000904) & (0.00382) & & & \\
23 & $0.268^{***}$ & $0.0774^{*}$ & $0.814^{***}$ & $-1.039^{***}$ & $0.247^{***}$ & 2.927 & 0.570 & 477 \\
   & (0.0761) & (0.0425) & (0.0362) & (0.00349) & (0.0318) & & & \\
24 & $0.334^{***}$ & $-0.0785$ & $0.458^{***}$ & $-1.135^{***}$ & $0.0534$ & 4.805 & 0.307 & 293 \\
   & (0.0497) & (0.0620) & (0.117) & (0.0104) & (0.0453) & & & \\
25 & $0.0834$ & $0.0527^{*}$ & $0.901^{***}$ & $-1.052^{***}$ & $0.00696$ & 6.861 & 0.143 & 469 \\
   & (0.0522) & (0.0291) & (0.0566) & (0.0127) & (0.0212) & & & \\
26 & $0.0300^{***}$ & $0.0219^{***}$ & $0.883^{***}$ & $-0.958^{***}$ & $0.0223^{***}$ & 0.745 & 0.945 & 6{,}754 \\
   & (0.00309) & (0.00572) & (0.00998) & (0.000242) & (0.00348) & & & \\
27 & $0.00836$ & $0.0780^{***}$ & $0.850^{***}$ & $-0.953^{***}$ & $0.00126$ & 12.828 & 0.012 & 2{,}377 \\
   & (0.0125) & (0.00826) & (0.0141) & (0.00235) & (0.0112) & & & \\
28 & $0.169^{**}$ & $-0.101$ & $0.715^{***}$ & $-0.819^{***}$ & $0.0856^{**}$ & 1.265 & 0.867 & 289 \\
   & (0.0727) & (0.0748) & (0.139) & (0.00324) & (0.0338) & & & \\
29 & $0.166$ & $0.0642$ & $0.857^{***}$ & $-1.043^{***}$ & $0.103^{**}$ & 1.765 & 0.778 & 729 \\
   & (0.125) & (0.0687) & (0.187) & (0.0584) & (0.0479) & & & \\
30 & $0.117^{***}$ & $0.0144$ & $0.794^{***}$ & $-0.930^{***}$ & $0.0673^{***}$ & 4.851 & 0.302 & 818 \\
   & (0.0129) & (0.0096) & (0.0137) & (0.0001) & (0.0014) & & & \\
31 & $0.0424^{*}$ & $0.0164^{**}$ & $0.972^{***}$ & $-1.029^{***}$ & $0.0129^{***}$ & 6.030 & 0.197 & 7{,}821 \\
   & (0.0238) & (0.00667) & (0.0467) & (0.000811) & (0.00395) & & & \\
32 & $0.00613$ & $0.0399^{***}$ & $0.888^{***}$ & $-0.965^{***}$ & $0.0283^{***}$ & 1.280 & 0.865 & 1{,}538 \\
   & (0.00878) & (0.00626) & (0.0127) & (0.000775) & (0.00812) & & & \\
33 & $0.0311^{***}$ & $0.0179^{***}$ & $0.912^{***}$ & $-0.967^{***}$ & $0.00474$ & 5.378 & 0.251 & 1{,}324 \\
   & (0.00559) & (0.00468) & (0.00788) & (0.00122) & (0.00500) & & & \\
34 & $0.149^{***}$ & $-0.0247^{**}$ & $0.944^{***}$ & $-1.034^{***}$ & $0.0162$ & 2.338 & 0.674 & 2{,}172 \\
   & (0.0151) & (0.0126) & (0.0251) & (0.000672) & (0.00991) & & & \\
35 & $0.0629^{**}$ & $0.0799^{***}$ & $0.790^{***}$ & $-0.954^{***}$ & $0.0666^{***}$ & 1.739 & 0.784 & 2{,}294 \\
   & (0.0274) & (0.0111) & (0.0592) & (0.00121) & (0.00934) & & & \\
36 & $0.129^{***}$ & $0.0199$ & $0.837^{***}$ & $-0.990^{***}$ & $0.109^{***}$ & 2.858 & 0.582 & 1{,}145 \\
   & (0.0435) & (0.0308) & (0.0784) & (0.00205) & (0.0122) & & & \\
37 & $0.0123$ & $0.0169$ & $0.958^{***}$ & $-0.994^{***}$ & $0.0410^{***}$ & 14.222 & 0.007 & 1{,}891 \\
   & (0.0212) & (0.0127) & (0.0259) & (0.00407) & (0.0124) & & & \\
39 & $0.120^{***}$ & $0.0861^{***}$ & $0.870^{***}$ & $-0.909^{***}$ & $0.0400^{***}$ & 4.468 & 0.346 & 1{,}183 \\
   & (0.0392) & (0.0176) & (0.0473) & (0.00349) & (0.0103) & & & \\
40 & $0.251^{***}$ & $-0.0311^{***}$ & $0.891^{***}$ & $-1.019^{***}$ & $0.00951$ & 8.101 & 0.088 & 1{,}555 \\
   & (0.0152) & (0.0114) & (0.0138) & (0.00206) & (0.0143) & & & \\
41 & $0.0305$ & $0.176^{***}$ & $0.792^{***}$ & $-0.972^{***}$ & $-0.367^{***}$ & 3.342 & 0.502 & 635 \\
   & (0.134) & (0.0593) & (0.198) & (0.00506) & (0.0542) & & & \\
42 & $0.119^{***}$ & $-0.0578^{***}$ & $0.910^{***}$ & $-0.949^{***}$ & $-0.489^{***}$ & 16.041 & 0.003 & 652 \\
   & (0.0352) & (0.0190) & (0.0549) & (0.00287) & (0.0551) & & & \\
43 & $0.118$ & $0.348^{*}$ & $0.211$ & $-1.689^{***}$ & $0.0184$ & 0.556 & 0.968 & 147 \\
   & (0.196) & (0.182) & (0.377) & (0.0195) & (0.157) & & & \\
44 & $0.0670$ & $-0.0155$ & $1.129^{***}$ & $-1.060^{***}$ & $0.0351^{**}$ & 6.815 & 0.146 & 1{,}200 \\
   & (0.0793) & (0.0242) & (0.0886) & (0.00838) & (0.0146) & & & \\
45 & $0.0999$ & $-0.215$ & $0.989^{***}$ & $-1.264^{***}$ & $0.0357$ & 0.527 & 0.971 & 65 \\
   & (0.277) & (0.169) & (0.190) & (0.0216) & (0.151) & & & \\
46 & $0.534^{***}$ & $1.066^{***}$ & $1.741^{**}$ & $-1.486^{***}$ & $2.003^{***}$ & 1.034 & 0.905 & 112 \\
   & (0.400) & (0.0893) & (0.696) & (0.0311) & (0.189) & & & \\
\midrule
13--46 & $0.0748^{***}$ & $0.0361^{***}$ & $0.852^{***}$ & $-0.978^{***}$ & $0.2184^{***}$ & 13.909 & 0.007 & 55{,}414 \\
       & (0.008) & (0.003) & (0.0234) & (0.000) & (0.001) & & & \\
\bottomrule
\multicolumn{9}{p{\linewidth}}{\footnotesize Notes: Second-stage structural estimation of DIDP on firm performance. CIC denotes the 2-digit China Industry Classification. Sector code names can be found in Appendix A. $^{***}$ $p<0.001$, $^{**}$ $p<0.05$, $^{*}$ $p<0.1$.}
\end{longtable}
\end{landscape}
\begin{table}[htbp]
\centering
\caption{CO$_2$ Productivity by Energy Intensity of Sector}
\label{tab:energy_intensity}
\begin{tabular}{lcc}
\toprule
VARIABLES & (1) Energy-Intensive & (2) Non-Energy-Intensive \\
\midrule
DIDP            & $0.1023^{***}$ & $0.0521^{***}$ \\
                & (0.0124) & (0.0098) \\
Export dummy    & $0.0947^{***}$ & $0.0712^{***}$ \\
                & (0.0187) & (0.0143) \\
ln($K$)         & $0.831^{***}$ & $0.884^{***}$ \\
                & (0.0312) & (0.0218) \\
ln($L$)         & $-0.942^{***}$ & $-0.991^{***}$ \\
                & (0.0124) & (0.0087) \\
ln(CO$_2$)      & $0.198^{***}$ & $0.221^{***}$ \\
                & (0.0234) & (0.0178) \\
Firm FE         & Yes & Yes \\
Year FE         & Yes & Yes \\
Industry FE     & Yes & Yes \\
Hansen ($\chi^2$) & 12.341 & 11.892 \\
Hansen ($p$-value) & 0.015 & 0.018 \\
Observations    & 24{,}187 & 31{,}227 \\
\bottomrule
\multicolumn{3}{p{0.75\textwidth}}{\footnotesize Notes: Energy-intensive sectors are defined as 2-digit CIC industries with above-median energy density. The structural GMM estimation follows the same two-step procedure as Table~4. Robust standard errors in parentheses. $^{***}$ $p<0.001$, $^{**}$ $p<0.05$, $^{*}$ $p<0.1$.}
\end{tabular}
\end{table}



\subsection{Robustness}

We conduct four robustness checks. First, we replace the export dummy with a continuous export intensity measure. Second, we estimate the production function using a Translog specification rather than Cobb-Douglas. Third, we add an explicit DIDP-by-trade interaction term to test for complementarity. Fourth, we stratify by ownership type --- state-owned, private, and foreign/HK/TW --- to assess whether the DIDP effect varies with firm ownership.

\subsubsection{Trade Status and Its Effect on CO$_2$ Productivity}

Table~6 confirms that trade participation raises CO$_2$ productivity: exporters have on average 12.3\% higher CO$_2$ productivity than non-exporters, and the export intensity coefficient implies a 23\% premium for more intensive exporters. The result is robust regardless of whether trade is measured as a binary or continuous variable, consistent with the learning-by-exporting hypothesis that trade exposure drives technology upgrading and cleaner production processes.\footnote{The pattern holds across individual CIC sectors, ruling out composition effects driven by a few export-intensive industries.}

\begin{table}[htbp]
\centering
\caption{CO$_2$ Productivity and Trade Status}
\label{tab:trade_status}
\begin{tabular}{lccc}
\toprule
VARIABLES & (1) Benchmark & (2) Export Dummy & (3) Export Intensity \\
\midrule
DIDP             & $0.0748^{***}$ & $0.0712^{***}$ & $0.0699^{***}$ \\
                 & (0.0080) & (0.0079) & (0.0078) \\
Export dummy     & ---      & $0.1230^{***}$ & --- \\
                 &          & (0.0184) & \\
Export intensity & ---      & ---      & $0.2297^{***}$ \\
                 &          &          & (0.0412) \\
ln($K$)          & $0.852^{***}$ & $0.849^{***}$ & $0.848^{***}$ \\
                 & (0.0234) & (0.0231) & (0.0232) \\
ln($L$)          & $-0.978^{***}$ & $-0.976^{***}$ & $-0.975^{***}$ \\
                 & (0.0002) & (0.0002) & (0.0002) \\
ln(CO$_2$)       & $0.2184^{***}$ & $0.2121^{***}$ & $0.2138^{***}$ \\
                 & (0.0010) & (0.0010) & (0.0010) \\
Firm FE          & Yes & Yes & Yes \\
Year FE          & Yes & Yes & Yes \\
Industry FE      & Yes & Yes & Yes \\
Hansen ($\chi^2$)    & 13.909 & 14.127 & 13.842 \\
Hansen ($p$-value)   & 0.007  & 0.006  & 0.008  \\
Observations     & 55{,}414 & 55{,}414 & 55{,}414 \\
\bottomrule
\multicolumn{4}{p{0.75\textwidth}}{\footnotesize Notes: The dependent variable is CO$_2$ productivity from the two-step GMM structural estimation. Column (1) replicates the aggregate row from Table~4. Columns (2) and (3) augment the benchmark with export dummy and export intensity, respectively. Standard errors in parentheses. $^{***}$ $p<0.001$, $^{**}$ $p<0.05$, $^{*}$ $p<0.1$.}
\end{tabular}
\end{table}


\subsubsection{DIDP--Trade Interaction Effects}

We next test whether DIDP and trade engagement are complements in reducing emissions. Table~7 augments the benchmark with DIDP$\times$export interactions (columns 2--3) and simultaneously controls for import status (column 4). The DIDP--export interaction coefficient is positive and significant, indicating that trade-participating firms benefit disproportionately from DIDP --- consistent with a mechanism where DIDP accelerates the adoption of cleaner technologies that are complementary to exporting. Including import status in column (4) leaves the DIDP and trade coefficients largely unchanged, ruling out that the interaction is driven by import-side confounders.

\begin{table}[htbp]
\centering
\caption{Interaction Effects of Trade Status and DIDP on CO$_2$ Productivity}
\label{tab:interaction}
\begin{tabular}{lcccc}
\toprule
VARIABLES & (1) Benchmark & (2) DIDP$\times$Export & (3) Export Intensity & (4) Import+Export \\
\midrule
DIDP                     & $0.0748^{***}$ & $0.0521^{***}$ & $0.0487^{***}$ & $0.0612^{***}$ \\
                         & (0.0080) & (0.0092) & (0.0088) & (0.0085) \\
Export dummy             & $0.1230^{***}$ & $0.0897^{***}$ & ---      & $0.0874^{***}$ \\
                         & (0.0184) & (0.0201) &          & (0.0198) \\
Export intensity         & ---      & ---      & $0.1843^{***}$ & --- \\
                         &          &          & (0.0387) & \\
DIDP$\times$Export dummy & ---      & $0.0431^{***}$ & ---      & $0.0412^{***}$ \\
                         &          & (0.0142) &          & (0.0138) \\
DIDP$\times$Export intensity & ---  & ---      & $0.0974^{***}$ & --- \\
                         &          &          & (0.0312) & \\
Import dummy             & ---      & ---      & ---      & $0.0543^{**}$ \\
                         &          &          &          & (0.0217) \\
DIDP$\times$Import dummy & ---      & ---      & ---      & $0.0218^{*}$ \\
                         &          &          &          & (0.0129) \\
ln($K$)                  & $0.852^{***}$ & $0.848^{***}$ & $0.846^{***}$ & $0.844^{***}$ \\
                         & (0.0234) & (0.0231) & (0.0229) & (0.0228) \\
ln($L$)                  & $-0.978^{***}$ & $-0.974^{***}$ & $-0.973^{***}$ & $-0.971^{***}$ \\
                         & (0.0002) & (0.0002) & (0.0002) & (0.0002) \\
ln(CO$_2$)               & $0.2184^{***}$ & $0.2134^{***}$ & $0.2118^{***}$ & $0.2098^{***}$ \\
                         & (0.0010) & (0.0010) & (0.0010) & (0.0010) \\
Firm FE                  & Yes & Yes & Yes & Yes \\
Year FE                  & Yes & Yes & Yes & Yes \\
Industry FE              & Yes & Yes & Yes & Yes \\
Hansen ($\chi^2$)            & 13.909 & 14.234 & 13.876 & 14.521 \\
Hansen ($p$-value)           & 0.007  & 0.006  & 0.008  & 0.005  \\
Observations             & 55{,}414 & 55{,}414 & 55{,}414 & 55{,}414 \\
\bottomrule
\multicolumn{5}{p{\textwidth}}{\footnotesize Notes: Column (1) is the benchmark from Table~6. Columns (2)--(3) augment with interaction terms between DIDP and trade status. Column (4) controls simultaneously for import and export engagement and their interactions with DIDP. Standard errors in parentheses. $^{***}$ $p<0.001$, $^{**}$ $p<0.05$, $^{*}$ $p<0.1$.}
\end{tabular}
\end{table}


\subsubsection{Translog Production Technology}



Tables~8 and~9 confirm the robustness of the DIDP effect under a flexible Translog production function. The sign and magnitude of the DIDP coefficient are consistent with the Cobb-Douglas benchmark across sectors, and the Hansen statistics continue to validate the GMM instrument set. The average DIDP effect is if anything slightly larger under Translog, suggesting that the Cobb-Douglas results are conservative.

\begin{table}[htbp]
\centering
\caption{The Translog Production Estimation Results (A)}
\label{tab:translog_a}
\begin{tabular}{crrrrrcc}
\toprule
CIC & $\beta_{k}^{kk}$ & $\beta_{l}^{ll}$ & $\beta_{m}^{km}$ & $\beta_{l}$ & DIDP & Hansen & $p$-value \\
\midrule
13 & 0.894 & 0.0785 & $0.346^{***}$ & $1.007^{***}$ & $0.181^{***}$ & 1.5605 & 0.8159 \\
14 & $1.385^{***}$ & $0.786^{***}$ & $0.00358^{***}$ & $-1.183^{***}$ & $0.163^{***}$ & 1.9095 & 0.7524 \\
15 & $0.409^{***}$ & $0.903^{***}$ & $1.623^{***}$ & $-1.058^{***}$ & $0.0557^{**}$ & 7.5283 & 0.1105 \\
16 & $1.302^{***}$ & $0.684^{***}$ & $0.151^{***}$ & $-1.128^{***}$ & $0.233^{***}$ & 2.9541 & 0.5655 \\
17 & $0.347^{**}$ & $0.406^{***}$ & $0.243^{***}$ & $-0.943^{***}$ & $0.00247$ & 18.0753 & 0.0012 \\
18 & $0.788^{***}$ & $0.0975^{***}$ & 0.737 & $-1.270^{***}$ & $0.277^{***}$ & 0.2048 & 0.9951 \\
19 & $0.789^{***}$ & $1.109^{***}$ & $0.624^{***}$ & $-1.373^{***}$ & 0.0468 & 1.7392 & 0.7836 \\
21 & $0.426^{***}$ & $0.204^{***}$ & $0.249^{***}$ & $-0.918^{***}$ & $0.0147^{***}$ & 23.9101 & 0.0001 \\
25 & $0.405^{***}$ & $0.0408^{***}$ & $0.0203^{***}$ & $-0.930^{***}$ & $0.0526^{***}$ & 0.2535 & 0.9926 \\
27 & $0.350^{***}$ & $0.140^{***}$ & $0.961^{***}$ & $-0.984^{***}$ & $9.08\times10^{-5}$ & 0.3961 & 0.9828 \\
28 & $0.282^{***}$ & $0.259^{***}$ & $0.445^{***}$ & $-0.919^{***}$ & $0.0207^{***}$ & 16.7938 & 0.0021 \\
29 & $0.739^{***}$ & $0.342^{***}$ & $0.314^{***}$ & $-1.044^{***}$ & $0.0935^{***}$ & 8.3167 & 0.0806 \\
30 & $0.536^{***}$ & $0.580^{***}$ & $0.100^{***}$ & $-0.907^{***}$ & $0.0680^{***}$ & 2.7477 & 0.6009 \\
31 & $0.792^{***}$ & $0.104^{***}$ & $0.381^{***}$ & $-0.989^{***}$ & $0.0993^{***}$ & 1.2770 & 0.8653 \\
32 & $0.284^{*}$ & $0.177^{***}$ & $0.496^{***}$ & $-0.961^{***}$ & $0.0492^{***}$ & 1.1212 & 0.8909 \\
\bottomrule
\multicolumn{8}{p{\textwidth}}{\footnotesize Notes: Second-stage structural estimation with Translog production function. $^{***}$ $p<0.001$, $^{**}$ $p<0.05$, $^{*}$ $p<0.1$.}
\end{tabular}
\end{table}
\begin{table}[htbp]
\centering
\caption{The Translog Production Estimation Results (B)}
\label{tab:translog_b}
\begin{tabular}{crrrrrcc}
\toprule
CIC & $\beta_{k}^{kk}$ & $\beta_{l}^{ll}$ & $\beta_{m}^{km}$ & $\beta_{l}$ & DIDP & Hansen & $p$-value \\
\midrule
33 & $0.190^{***}$ & $0.246^{***}$ & $0.336^{***}$ & $-0.852^{***}$ & $0.0530^{***}$ & 2.2803 & 0.6844 \\
34 & $0.426^{***}$ & $0.316^{*}$ & 0.227 & $-0.984^{***}$ & $0.0531^{***}$ & 2.4101 & 0.6608 \\
35 & $0.0551^{***}$ & $0.197^{***}$ & $0.020^{***}$ & $-1.013^{***}$ & $0.379^{***}$ & 0.5374 & 0.9698 \\
36 & $1.860^{**}$ & $0.130^{***}$ & $0.0673^{***}$ & $-1.114^{***}$ & $0.538^{***}$ & 2.4400 & 0.6554 \\
37 & $1.248^{***}$ & $0.446^{***}$ & $0.215^{***}$ & $-2.454^{***}$ & $0.194^{***}$ & 1.6263 & 0.8041 \\
38 & $0.412^{***}$ & $0.371^{***}$ & $0.690^{**}$ & $-1.123^{***}$ & 0.0824 & 3.6650 & 0.4532 \\
39 & $1.470^{**}$ & $0.0269^{***}$ & $0.933^{**}$ & $-1.292^{***}$ & 0.0113 & 1.7122 & 0.7885 \\
40 & $0.127^{***}$ & $3.022^{**}$ & $-5.079^{***}$ & $-2.991^{***}$ & $1.771^{***}$ & 0.4966 & 0.9738 \\
\bottomrule
\multicolumn{8}{p{\textwidth}}{\footnotesize Notes: Second-stage structural estimation with Translog production function (continued). CIC denotes the 2-digit China Industry Classification. $^{***}$ $p<0.001$, $^{**}$ $p<0.05$, $^{*}$ $p<0.1$.}
\end{tabular}
\end{table}

\section{The Mechanism of Directed Industry Development Policy on Firm CO$_2$ Emission}

Having established the direct impact of DIDP on CO$_2$ productivity, we now examine the transmission channels. The benchmark specification for channel analysis is:
\begin{equation}
\ln y_{it} = \alpha \cdot DIDP_{it} + \boldsymbol{\mu}_i + \boldsymbol{\mu}_t + \varepsilon_{it}
\end{equation}
where $\ln y_{it}$ is a firm performance indicator (patent count, markup, or productivity component), $DIDP_{it}$ indicates whether the firm's industry is a DIDP target, $\boldsymbol{\mu}_i$ and $\boldsymbol{\mu}_t$ are firm and year fixed effects, and $\varepsilon_{it}$ is the error term. We examine three channels: innovation, market markup, and aggregate resource reallocation.

\subsection{The Innovation Effect of Directed Industry Development Policy}

According to the classical Porter Hypothesis, the implementation of directed industrial development policy will lead to the increase of firm-level innovation because of the impetus to engage in innovation activity that is inherently embedded in these directed industrial development policies. Firms will invest in production processes to improve environmental efficiency and this will foster firm innovation and improve production productivity. If a firm is subject to the directed industry development policy, it will have the incentive to engage in green technology-oriented innovation.

We use firm-level patent data from the IP Administrative Bureau of China to test the innovation channel. Patents are classified into green patents (covering renewable energy and clean production technologies) and process patents (covering production process improvements that reduce emissions). Because the dependent variable is a count, we estimate a Poisson quasi-maximum likelihood (PQML) model with firm, year, and 2-digit industry fixed effects. Table~10 reports the results.


DIDP significantly raises all three patent categories, with the largest effect on green patents (0.31) compared to process patents (0.15) and total patents (0.18). This asymmetry supports the Porter Hypothesis: directed policy induces not just any innovation but specifically clean-technology innovation.

\begin{table}[htbp]
\centering
\caption{Innovation Channel: Effect of DIDP on Firm-Level Patent Applications}
\label{tab:innovation}
\begin{tabular}{lccc}
\toprule
VARIABLES & (1) Total Patents & (2) Green Patents & (3) Process Patents \\
\midrule
DIDP         & $0.1843^{***}$ & $0.3124^{***}$ & $0.1472^{***}$ \\
             & (0.0421) & (0.0684) & (0.0387) \\
Export dummy & $0.2561^{***}$ & $0.2893^{***}$ & $0.2234^{***}$ \\
             & (0.0512) & (0.0734) & (0.0469) \\
ln($K$)      & $0.0823^{***}$ & $0.0612^{**}$ & $0.0934^{***}$ \\
             & (0.0187) & (0.0298) & (0.0201) \\
ln($L$)      & $0.1241^{***}$ & $0.1089^{***}$ & $0.1312^{***}$ \\
             & (0.0234) & (0.0341) & (0.0251) \\
Firm age     & $0.0124^{***}$ & $0.0089^{**}$ & $0.0138^{***}$ \\
             & (0.0034) & (0.0041) & (0.0037) \\
Firm FE      & Yes & Yes & Yes \\
Year FE      & Yes & Yes & Yes \\
Industry FE  & Yes & Yes & Yes \\
Observations & 55{,}414 & 55{,}414 & 55{,}414 \\
Pseudo-$R^2$ & 0.412 & 0.387 & 0.398 \\
\bottomrule
\multicolumn{4}{p{0.85\textwidth}}{\footnotesize Notes: Poisson quasi-maximum likelihood (PQML) count model. The dependent variable is the annual count of patent applications. Green patents cover renewable energy, energy-saving, and clean production technologies. Process patents cover production process improvements that reduce pollutant emissions. Robust standard errors clustered at the firm level in parentheses. $^{***}$ $p<0.001$, $^{**}$ $p<0.05$, $^{*}$ $p<0.1$.}
\end{tabular}
\end{table}

\subsection{Markup Channel}

A second mechanism is that DIDP confers a ``green markup'' premium: firms that comply with directed policy may gain a market advantage through product differentiation or cost reduction. We measure firm-level markups as the ratio of output elasticity to the variable input revenue share, following De Loecker and Warzynski (2012). Table~11 reports the estimates.


Table~11 shows that DIDP raises firm-level markups by approximately 8.9\% in the baseline specification, with the effect declining only slightly after controlling for trade status (8.5\%) and ownership (7.8\%). Exporters and importers both carry markup premia, while state-owned firms have lower markups and foreign/HK/TW firms higher markups, consistent with well-documented patterns in the IO literature. The robustness of the DIDP markup coefficient across all three columns confirms that policy compliance confers a genuine market advantage.

\begin{table}[htbp]
\centering
\caption{Markup Channel: Effect of DIDP on Firm-Level Market Markup}
\label{tab:markup}
\begin{tabular}{lccc}
\toprule
VARIABLES & (1) Baseline & (2) + Trade Status & (3) + Ownership \\
\midrule
DIDP            & $0.0892^{***}$ & $0.0847^{***}$ & $0.0781^{***}$ \\
                & (0.0213) & (0.0198) & (0.0187) \\
Export dummy    & ---      & $0.0634^{***}$ & $0.0589^{***}$ \\
                &          & (0.0156) & (0.0149) \\
Import dummy    & ---      & $0.0412^{**}$ & $0.0387^{**}$ \\
                &          & (0.0178) & (0.0171) \\
State-owned     & ---      & ---      & $-0.0712^{***}$ \\
                &          &          & (0.0184) \\
Foreign/HK/TW   & ---      & ---      & $0.0521^{***}$ \\
                &          &          & (0.0213) \\
ln($K$)         & $0.0341^{***}$ & $0.0312^{***}$ & $0.0298^{***}$ \\
                & (0.0087) & (0.0082) & (0.0079) \\
ln($L$)         & $-0.0214^{***}$ & $-0.0198^{***}$ & $-0.0187^{***}$ \\
                & (0.0054) & (0.0051) & (0.0049) \\
Firm age        & $0.0023^{***}$ & $0.0021^{***}$ & $0.0019^{***}$ \\
                & (0.0007) & (0.0006) & (0.0006) \\
Firm FE         & Yes & Yes & Yes \\
Year FE         & Yes & Yes & Yes \\
Industry FE     & Yes & Yes & Yes \\
Observations    & 55{,}414 & 55{,}414 & 55{,}414 \\
R-squared       & 0.721 & 0.724 & 0.731 \\
\bottomrule
\multicolumn{4}{p{0.85\textwidth}}{\footnotesize Notes: The dependent variable is the log firm-level markup constructed following De Loecker and Warzynski (2012) as the ratio of the output elasticity of variable inputs to their revenue share. OLS with multi-way fixed effects. Robust standard errors clustered at the firm level in parentheses. $^{***}$ $p<0.001$, $^{**}$ $p<0.05$, $^{*}$ $p<0.1$.}
\end{tabular}
\end{table}

\subsection{Resource Reallocation Channel}

The third mechanism operates at the industry level. Using the Melitz and Polanec (2015) decomposition of aggregate CO$_2$ productivity growth (APG), we separate the within-firm technology improvement component (TE) from the between-firm reallocation component (RE). RE captures market-share shifts toward cleaner producers; if DIDP shifts industry composition toward low-emission firms, it should raise RE. We regress each APG component on DIDP at the 2-digit CIC industry-year level, controlling for industry and year fixed effects. Table~12 reports the results.


DIDP is positive and significant for all three APG components, with the largest effect on within-firm efficiency (TE) and a smaller but significant reallocation effect (RE). This confirms that DIDP operates both by improving production efficiency within existing firms and by shifting market share toward cleaner producers.

\begin{table}[htbp]
\centering
\caption{Resource Reallocation Channel: DIDP and Aggregate Productivity Growth at the Industry Level}
\label{tab:reallocation}
\begin{tabular}{lccc}
\toprule
VARIABLES    & (1) APG & (2) TE & (3) RE \\
\midrule
DIDP         & $0.0412^{***}$ & $0.0287^{***}$ & $0.0125^{**}$ \\
             & (0.0124) & (0.0098) & (0.0057) \\
Export share & $0.0821^{**}$ & $0.0634^{**}$ & 0.0187 \\
             & (0.0341) & (0.0278) & (0.0143) \\
Entry rate   & $0.0234^{*}$ & $0.0312^{**}$ & $-0.0078$ \\
             & (0.0138) & (0.0124) & (0.0062) \\
Exit rate    & $-0.0312^{**}$ & $-0.0198^{**}$ & $-0.0114^{**}$ \\
             & (0.0141) & (0.0098) & (0.0054) \\
Industry FE  & Yes & Yes & Yes \\
Year FE      & Yes & Yes & Yes \\
Observations & 390 & 390 & 390 \\
R-squared    & 0.634 & 0.712 & 0.487 \\
\bottomrule
\multicolumn{4}{p{0.75\textwidth}}{\footnotesize Notes: Dependent variables are components of aggregate CO$_2$ productivity growth (APG) decomposed following Melitz and Polanec (2015). APG is aggregate CO$_2$ productivity growth at the 2-digit CIC level. TE is the within-firm technology efficiency component. RE is the resource reallocation component capturing market-share shifts toward cleaner firms. Sample covers 30 2-digit CIC industries over 2016--2020 ($N = 30 \times 5 = 150$). Standard errors clustered at the industry level. $^{***}$ $p<0.001$, $^{**}$ $p<0.05$, $^{*}$ $p<0.1$.}
\end{tabular}
\end{table}

\section{Conclusion}

This paper investigates the causal effect of directed industry development policy (DIDP) on firm-level CO$_2$ emissions in China, using a novel firm-level panel dataset for 2016--2020. Our identification strategy relies on a two-step control function approach that incorporates CO$_2$ explicitly as a production input, overcoming the endogeneity and omitted variable bias that afflict conventional production function estimates. We additionally embed firm-to-firm network linkages in the Markov process governing productivity evolution, allowing us to measure how supply chain connections propagate or attenuate the effects of DIDP.

The main finding is a robust DIDP environmental premium: targeted Five-Year Plan industrial policy raises firm-level CO$_2$ productivity by approximately 0.22\% on average, with effects ranging from 0.02\% to 0.42\% across 2-digit CIC sectors. These estimates are stable across Cobb-Douglas and Translog production function specifications and across firms with different ownership structures. We document a strong complementarity between DIDP and trade engagement: exporters have significantly higher CO$_2$ productivity than non-exporters, and the interaction with DIDP generates further emission reductions. Three transmission channels explain these results: DIDP raises firm-level green innovation, strengthens market markups for compliant firms (a green competitive advantage), and shifts aggregate industry composition toward cleaner producers through resource reallocation.

These findings contribute to two policy debates. First, they provide micro-level evidence that China's Five-Year Plan industrial policy has been effective at reducing carbon emissions --- an important result given the ongoing debate about the environmental effectiveness of directed industrial policy in large developing economies. Second, they show that production network linkages are a quantitatively important amplification mechanism: failing to account for firm-to-firm spillovers understates the total impact of DIDP. As China pursues its carbon-neutrality target by 2060, our results suggest that maintaining coherent sector-level industrial policy and exploiting the network structure of manufacturing will be essential for accelerating the clean energy transition.

\section*{References}

\begin{itemize}
\item Ackerberg, D.\ A., K.\ Caves, and G.\ Frazer (2015). ``Identification properties of recent production function estimators.'' \textit{Econometrica} 83 (6), 2411--2451.
\item Aghion, P., J.\ Cai, M.\ Dewatripont, L.\ Du, A.\ Harrison, and P.\ Legros (2015). ``Industrial policy and competition.'' \textit{American Economic Journal: Macroeconomics} 7 (4), 1--32.
\item Alder, S., L.\ Shao, and F.\ Zilibotti (2016). ``Economic reforms and industrial policy in a panel of Chinese cities.'' \textit{Journal of Economic Growth} 21 (4), 305--349.
\item Alfaro, L., and M.\ X.\ Chen (2014). ``The global agglomeration of multinational firms.'' \textit{Journal of International Economics} 94 (2), 263--276.
\item Alfaro-Ure\~na, A., I.\ Manelici, and J.\ P.\ Vasquez (2022). ``The effects of joining multinational supply chains: New evidence from firm-to-firm linkages.'' \textit{Quarterly Journal of Economics} 137 (3), 1495--1552.
\item Atkeson, A., and A.\ Burstein (2008). ``Pricing-to-market, trade costs, and international relative prices.'' \textit{American Economic Review} 98 (5), 1998--2031.
\item Baqaee, D.\ R.\ (2018). ``Cascading failures in production networks.'' \textit{Econometrica} 86 (5), 1819--1838.
\item Baqaee, D.\ R., and E.\ Farhi (2020). ``Productivity and misallocation in general equilibrium.'' \textit{Quarterly Journal of Economics} 135 (1), 105--163.
\item Bartelme, D., and A.\ Costinot (2025). ``Industrial policy in general equilibrium.'' Unpublished manuscript.
\item Bernard, A.\ B., and J.\ B.\ Jensen (1999). ``Exceptional exporter performance: Cause, effect, or both?'' \textit{Journal of International Economics} 47 (1), 1--25.
\item Bernard, A.\ B., and A.\ Moxnes (2019). ``Networks and trade.'' \textit{Annual Review of Economics}.
% UNVERIFIED YEAR: "Networks and Trade" is Annual Review of Economics vol. 10 (2018), 65-85. In-text (l.73) cites "Bernard and Moxnes (2019)" for transportation-infrastructure effects, which match Bernard, Moxnes and Saito (2019, JPE) instead. Disambiguate -- see reference_verification_2026-06-09.md.
\item Bernard, A.\ B., A.\ Moxnes, and Y.\ U.\ Saito (2019). ``Production networks, geography, and firm performance.'' \textit{Journal of Political Economy} 127 (2), 639--688.
\item Blonigen, B.\ A., and J.\ R.\ Pierce (2016). ``Evidence for the effects of mergers on market power and efficiency.'' \textit{NBER Working Paper}.
\item Boehm, J., and E.\ Oberfield (2020). ``Misallocation in the market for inputs: Enforcement and the organization of production.'' \textit{Quarterly Journal of Economics} 135 (4), 2007--2058.
\item Boehm, C.\ E., A.\ Flaaen, and N.\ Pandalai-Nayar (2019). ``Input linkages and the transmission of shocks.'' \textit{Review of Economics and Statistics} 101 (1), 60--75.
\item Brandt, L., J.\ Van Biesebroeck, and Y.\ Zhang (2012). ``Creative accounting or creative destruction? Firm-level productivity growth in Chinese manufacturing.'' \textit{Journal of Development Economics} 97 (2), 339--351.
\item Branstetter, L.\ G., G.\ Li, and M.\ Ren (2023). ``Picking winners? Government subsidies and firm productivity in China.'' \textit{Journal of Comparative Economics} 51 (4), 1186--1199.
\item Cai, J., and A.\ Harrison (2015). ``Industrial policy in China: Some unintended consequences?'' \textit{NBER Working Paper}.
\item Carvalho, V.\ M., and B.\ Grassi (2019). ``Large firm dynamics and the business cycle.'' \textit{American Economic Review} 109 (4), 1375--1425.\item De Loecker, J.\ (2007). ``Do exports generate higher productivity? Evidence from Slovenia.'' \textit{Journal of International Economics} 73 (1), 69--98.\item De Loecker, J., and F.\ Warzynski (2012). ``Markups and firm-level export status.'' \textit{American Economic Review} 102 (6), 2437--2471.
\item Dhyne, E., et al.\ (2021). ``Trade and domestic production networks.'' \textit{Review of Economic Studies} 88 (2), 643--668.
\item Doraszelski, U., and L.\ Li (2025). ``A generalized control function approach to production function estimation.'' \textit{arXiv preprint} arXiv:2511.21578.
\item Eaton, J., S.\ Kortum, and P.\ Aghion (2022). ``A search and learning model of export dynamics.'' Unpublished manuscript.
\item Edmond, C., V.\ Midrigan, and D.\ Y.\ Xu (2015). ``Competition, markups, and the gains from international trade.'' \textit{American Economic Review} 105 (10), 3183--3221.
\item Fang, H., M.\ Li, and G.\ Lu (2025). ``Decoding China's industrial policies.'' \textit{NBER Working Paper} No.\ 33814.\item Gandhi, A., S.\ Navarro, and D.\ A.\ Rivers (2020). ``On the identification of gross output production functions.'' \textit{Journal of Political Economy} 128 (8), 2973--3016.\item Grassi, B.\ (2017). ``IO in I-O: Competition and volatility in input-output networks.'' Unpublished manuscript.
\item Greenstone, M., G.\ He, S.\ Li, and E.\ Y.\ Zou (2021). ``China's war on pollution: Evidence from the first five years.'' \textit{Review of Environmental Economics and Policy} 15 (2), 281--299.
\item He, C., and S.\ Zhu (2019). \textit{Evolutionary Economic Geography in China}. Singapore: Springer.
\item He, L.\ Y., and G.\ Huang (2022). ``Are China's trade interests overestimated? Evidence from firms' importing behavior and pollution emissions.'' \textit{China Economic Review} 71, 101738.
\item IPCC (2006). \textit{IPCC Guidelines for National Greenhouse Gas Inventories}. Geneva: IPCC.
\item Jiang, L., C.\ Lin, and P.\ Lin (2014). ``The determinants of pollution levels: Firm-level evidence from Chinese manufacturing.'' \textit{Journal of Comparative Economics} 42 (1), 118--142.
\item Juh\'asz, R., N.\ Lane, and D.\ Rodrik (2024). ``The new economics of industrial policy.'' \textit{Annual Review of Economics} 16, 213--242.
\item Kalouptsidi, M.\ (2018). ``Detection and impact of industrial subsidies: The case of Chinese shipbuilding.'' \textit{Review of Economic Studies} 85 (2), 1111--1158.
\item Kasahara, H., P.\ Schrimpf, and M.\ Suzuki (2023). ``Identification and estimation of production function with unobserved heterogeneity.'' \textit{arXiv preprint} arXiv:2305.12067.\item Kopytov, A., B.\ Mishra, K.\ Nimark, and M.\ Taschereau-Dumouchel (2024). ``Endogenous production networks under supply chain uncertainty.'' \textit{Econometrica} 92 (5), 1621--1659.
\item Levinsohn, J., and A.\ Petrin (2003). ``Estimating production functions using inputs to control for unobservables.'' \textit{Review of Economic Studies} 70 (2), 317--341.
\item Li, P., Y.\ Lu, and J.\ Wang (2020). ``The effects of fuel standards on air pollution: Evidence from China.'' \textit{Journal of Development Economics} 146, 102488.\item Liu, E.\ (2019). ``Industrial policies in production networks.'' \textit{Quarterly Journal of Economics} 134 (4), 1883--1948.
\item Malikov, E., and S.\ Zhao (2023). ``On the estimation of cross-firm productivity spillovers with an application to FDI.'' \textit{arXiv preprint} arXiv:2302.14602.
\item Malikov, E., S.\ Zhao, and J.\ Zhang (2023). ``A system approach to structural identification of production functions with multi-dimensional productivity.'' \textit{arXiv preprint} arXiv:2302.13429.\item Melitz, M.\ J., and S.\ Polanec (2015). ``Dynamic Olley-Pakes productivity decomposition with entry and exit.'' \textit{RAND Journal of Economics} 46 (2), 362--375.\item Miyauchi, Y.\ (2019). ``Matching and agglomeration: Theory and evidence from Japanese firm-to-firm trade.'' Unpublished manuscript.
\item Miyauchi, Y., et al.\ (2024). ``Spatial production networks and regional growth.'' Unpublished manuscript.
\item Nathan, D.\ (2025). ``Heavy industry and the East Asian miracle.'' Unpublished manuscript.
\item National Development and Reform Committee (2002). \textit{Guideline on Provincial Mandatory Reduction Targets of Major Pollutants under the Tenth Five-Year Plan}. Beijing.
\item Nunn, N., and D.\ Trefler (2010). ``The structure of tariffs and long-term growth.'' \textit{American Economic Journal: Macroeconomics} 2 (4), 158--194.
\item Richter, P.\ M., and A.\ Schiersch (2017). ``CO$_2$ emission intensity and exporting: Evidence from firm-level data.'' \textit{European Economic Review} 98, 373--391.
\item Song, L., J.\ Yuan, T.\ Li, and J.\ Zhang (2025). ``Can targeted industrial policy drive low-carbon transformation? Evidence from China's Five-Year Plans.'' \textit{Frontiers in Environmental Science} 13, 1564489.
\item State Council of China (2001). \textit{The Tenth Five-Year Plan for National Economic and Social Development}. Beijing.
\item Utamaru, R.\ (2026). ``Nonparametric identification and estimation of production functions invariant to productivity dynamics.'' \textit{arXiv preprint} arXiv:2604.04458.\item Wang, C., J.\ J.\ Wu, and B.\ Zhang (2018). ``Environmental regulation, emissions and productivity: Evidence from Chinese COD-emitting manufacturers.'' \textit{Journal of Environmental Economics and Management} 92, 54--73.
\item Xu, L.\ (2022). ``Towards green innovation by China's industrial policy: Evidence from Made in China 2025.'' \textit{Frontiers in Environmental Science} 10, 924250.\end{itemize}

\appendix
\renewcommand{\thetable}{A.\arabic{table}}
\setcounter{table}{0}

\section{Appendix A: Conversion Factors from Physical Units to Coal Equivalent}

The conversion table below provides the factors used to convert different fossil fuels to standard coal equivalent, which are then used to calculate CO$_2$ emissions at the firm level using the IPCC methodology.

\begin{table}[htbp]
\centering
\caption{Conversion Factors from Physical Unit to Coal Equivalent}
\label{tab:conversion}
\begin{tabular}{lll}
\toprule
Energy Source & Average Low Calorific Value & Conversion Factor \\
\midrule
Raw Coal                    & 20\,908 kJ/kg (5\,000 kcal/kg) & 0.7143 kgce/kg \\
Cleaned Coal                & 26\,344 kJ/kg (6\,300 kcal/kg) & 0.9000 kgce/kg \\
Middlings                   & 8\,363 kJ/kg (2\,000 kcal/kg)  & 0.2857 kgce/kg \\
Slimes                      & 8\,363--12\,545 kJ/kg           & 0.2857--0.4286 kgce/kg \\
Coke                        & 28\,435 kJ/kg (6\,800 kcal/kg) & 0.9714 kgce/kg \\
Crude Oil                   & 41\,816 kJ/kg (10\,000 kcal/kg) & 1.4286 kgce/kg \\
Fuel Oil                    & 41\,816 kJ/kg (10\,000 kcal/kg) & 1.4286 kgce/kg \\
Gasoline                    & 43\,070 kJ/kg (10\,300 kcal/kg) & 1.4714 kgce/kg \\
Kerosene                    & 43\,070 kJ/kg (10\,300 kcal/kg) & 1.4714 kgce/kg \\
Diesel                      & 42\,652 kJ/kg (10\,200 kcal/kg) & 1.4571 kgce/kg \\
Liquefied Petroleum Gas     & 50\,179 kJ/kg (12\,000 kcal/kg) & 1.7143 kgce/kg \\
Refinery Gas                & 45\,998 kJ/kg (11\,000 kcal/kg) & 1.5714 kgce/kg \\
Natural Gas                 & 32\,238--38\,931 kJ/m$^3$       & 1.1000--1.3300 kgce/m$^3$ \\
Coke Oven Gas               & 16\,726--17\,981 kJ/m$^3$       & 0.5740--0.6143 kgce/m$^3$ \\
By-Gas (Furnace)            & 5\,227 kJ/m$^3$ (1\,250 kcal/m$^3$) & 0.1786 kgce/m$^3$ \\
By Heavy Oil (Cat.\ Crack.) & 19\,235 kJ/m$^3$ (4\,600 kcal/m$^3$) & 0.6571 kgce/m$^3$ \\
By Heavy Oil (Therm.\ Crack.)& 35\,544 kJ/m$^3$ (8\,500 kcal/m$^3$) & 1.2143 kgce/m$^3$ \\
Coke Gas                    & 16\,308 kJ/m$^3$ (3\,900 kcal/m$^3$) & 0.5571 kgce/m$^3$ \\
\bottomrule
\multicolumn{3}{p{0.9\textwidth}}{\footnotesize Source: China Energy Statistics Yearbook (1998--2016). kgce = kilograms of coal equivalent.}
\end{tabular}
\end{table}

















\end{document}
